\ifdefined\gkver\else\def\gkver{student}\fi
\documentclass[\gkver]{gaokaozhenti}
\begin{document}

\chapter{2023年重庆}

\begin{enumerate}
\item
%% number: 1
%% paperName: 2023年重庆高考第 1 题 4 分
%% typeId: 1
%% type: 单选题
%% chapter: 力物体的平衡
%% point: 力的合成
%% method: 无
%% score: 4
%% degree: 850
%% duplicateId: 0
%% body:
矫正牙齿时，可用牵引线对牙施加力的作用。若某颗牙齿受到牵引线的两个作用力大小均为 $F$，夹角为 $\alpha$（如图），则该牙所受两牵引力的合力大小为\xzanswer{B}

\onepicture[h=3.40cm]{figs/01.png}

\fourchoices[answer=B]
{$2F\sin\frac{\alpha}{2}$}
{$2F\cos\frac{\alpha}{2}$}
{$F\sin\alpha$}
{$F\cos\alpha$}

%% answer: B
%% memo:
\memoanswer{
根据平行四边形定则可知，该牙所受两牵引力的合力大小为
$F_{\text{合}}=2F\cos\frac{\alpha}{2}$

故选 B。
}

\item
%% number: 2
%% paperName: 2023年重庆高考第 2 题 4 分
%% typeId: 1
%% type: 单选题
%% chapter: 电磁感应
%% point: 法拉第电磁感应定律
%% method: 无
%% score: 4
%% degree: 750
%% duplicateId: 0
%% body:
某小组设计了一种呼吸监测方案：在人身上缠绕弹性金属线圈，观察人呼吸时处于匀强磁场中的线圈面积变化产生的电压，了解人的呼吸状况。如图所示，线圈 P 的匝数为 $N$，磁场的磁感应强度大小为 $B$，方向与线圈轴线的夹角为 $\theta$。若某次吸气时，在 $t$ 时间内每匝线圈面积增加了 $S$，则线圈 P 在该时间内的平均感应电动势为\xzanswer{A}

\onepicture[h=3.60cm]{figs/02.png}

\fourchoices[answer=A]
{$\frac{NBS\cos\theta}{t}$}
{$\frac{NBS\sin\theta}{t}$}
{$\frac{BS\sin\theta}{t}$}
{$\frac{BS\cos\theta}{t}$}

%% answer: A
%% memo:
\memoanswer{
根据法拉第电磁感应定律有
$\overline{E}=N\frac{\Delta\Phi}{\Delta t}=NB\cos\theta\frac{S}{t}=\frac{NBS\cos\theta}{t}$

故选 A。
}

\item
%% number: 3
%% paperName: 2023年重庆高考第 3 题 4 分
%% typeId: 1
%% type: 单选题
%% chapter: 电场
%% point: 电场叠加
%% method: 无
%% score: 4
%% degree: 650
%% duplicateId: 0
%% body:
真空中固定有两个点电荷，负电荷 $Q_{1}$ 位于坐标原点处，正电荷 $Q_{2}$ 位于 $x$ 轴上，$Q_{2}$ 的电荷量大小为 $Q_{1}$ 的 8 倍。若这两点电荷在 $x$ 轴正半轴的 $x=x_{0}$ 处产生的合电场强度为 0，则 $Q_{1}$、$Q_{2}$ 相距\xzanswer{B}

\fourchoices[answer=B]
{$\sqrt{2}x_{0}$}
{$(2\sqrt{2}-1)x_{0}$}
{$2\sqrt{2}x_{0}$}
{$(2\sqrt{2}+1)x_{0}$}

%% answer: B
%% memo:
\memoanswer{
依题意，两点电荷电性相反，且 $Q_{2}$ 的电荷量较大，所以合场强为 0 的位置应该在 $x$ 轴的负半轴，设两点电荷相距 $L$，根据点电荷场强公式可得
$\frac{kQ_{1}}{x_{0}^{2}}=\frac{kQ_{2}}{(x_{0}+L)^{2}}$

又 $Q_{2}=8Q_{1}$

解得 $L=(2\sqrt{2}-1)x_{0}$

故选 B。
}

\item
%% number: 4
%% paperName: 2023年重庆高考第 4 题 4 分
%% typeId: 1
%% type: 单选题
%% chapter: 气体性质
%% point: 气体图象
%% method: 无
%% score: 4
%% degree: 650
%% duplicateId: 0
%% body:
密封于气缸中的理想气体，从状态 a 依次经过 $ab$、$bc$ 和 $cd$ 三个热力学过程达到状态 d。若该气体的体积 $V$ 随热力学温度 $T$ 变化的 $V-T$ 图像如图所示，则对应的气体压强 $p$ 随 $T$ 变化的 $p-T$ 图像正确的是\xzanswer{C}

\onepicture[h=3.00cm]{figs/04.png}

\fourchoices[answer=C, ispicture=true, h=1.90cm]
{figs/04a.png}
{figs/04b.png}
{figs/04c.png}
{figs/04d.png}

%% answer: C
%% memo:
\memoanswer{
由 $V-T$ 图像可知，理想气体 $ab$ 过程做等压变化，$bc$ 过程做等温变化，$cd$ 过程做等容变化。根据理想气体状态方程，有
$\frac{pV}{T}=C$

可知 $bc$ 过程理想气体的体积增大，则压强减小。

故选 C。
}

\item
%% number: 5
%% paperName: 2023年重庆高考第 5 题 4 分
%% typeId: 1
%% type: 单选题
%% chapter: 光学
%% point: 全反射
%% method: 无
%% score: 4
%% degree: 550
%% duplicateId: 0
%% body:
某实验小组利用双缝干涉实验装置分别观察 a、b 两单色光的干涉条纹，发现在相同的条件下光屏上 a 光相邻两亮条纹的间距比 b 光的小。他们又将 a、b 光以相同的入射角由水斜射入空气，发现 a 光的折射角比 b 光的大，则\xzanswer{D}

\fourchoices[answer=D]
{在空气中传播时，a 光的波长比 b 光的大}
{在水中传播时，a 光的速度比 b 光的大}
{在水中传播时，a 光的频率比 b 光的小}
{由水射向空气时，a 光的全反射临界角比 b 光的小}

%% answer: D
%% memo:
\memoanswer{
A．根据相邻两条亮条纹的间距计算公式
$\Delta x=\frac{L}{d}\lambda$

由此可知
$\lambda_{a}<\lambda_{b}$

故 A 错误；

B．根据折射定律
$n=\frac{\sin i}{\sin r}$

a、b 光以相同的入射角由水斜射入空气，a 光的折射角比 b 光的大，则
$n_{a}>n_{b}$

根据光在介质中的传播速度与折射率的关系
$n=\frac{c}{v}$

可得在水中传播时，a 光的速度比 b 光的小，故 B 错误；

C．在水中传播时，a 光的折射率比 b 光的大，所以 a 光的频率比 b 光的大，故 C 错误；

D．根据临界角与折射率的关系
$n=\frac{1}{\sin C}$

可得在水中传播时，a 光的折射率比 b 光的大，a 光的全反射临界角比 b 光的小，故 D 正确。

故选 D。
}

\item
%% number: 6
%% paperName: 2023年重庆高考第 6 题 4 分
%% typeId: 1
%% type: 单选题
%% chapter: 原子和原子核
%% point: 天然放射性
%% method: 无
%% score: 4
%% degree: 550
%% duplicateId: 0
%% body:
原子核 \ce{^{235}_{92} U} 可以经过多次 $\alpha$ 和 $\beta$ 衰变成为稳定的原子核 \ce{^{207}_{82} Pb}，在该过程中，可能发生的 $\beta$ 衰变是\xzanswer{A}

\fourchoices[answer=A]
{\ce{^{223}_{87} Fr -> ^{223}_{88} Ra + ^{0}_{-1} e}}
{\ce{^{213}_{84} Bi -> ^{213}_{83} Po + ^{0}_{-1} e}}
{\ce{^{225}_{88} Ra -> ^{225}_{89} Ac + ^{0}_{-1} e}}
{\ce{^{218}_{84} Po -> ^{218}_{85} At + ^{0}_{-1} e}}

%% answer: A
%% memo:
\memoanswer{
原子核 \ce{^{235}_{92} U} 衰变成为稳定的原子核 \ce{^{207}_{82} Pb} 质量数减小了 28，则经过了 7 次 $\alpha$ 衰变，中间生成的新核的质量数可能为 231，227，223，219，215，211，则发生 $\beta$ 衰变的原子核的质量数为上述各数，则 BCD 都不可能，根据核反应的质量数和电荷数守恒可知，选项 A 反应正确。

故选 A。
}

\item
%% number: 7
%% paperName: 2023年重庆高考第 7 题 4 分
%% typeId: 1
%% type: 单选题
%% chapter: 电磁感应
%% point: 电磁感应中的变加速
%% method: 无
%% score: 4
%% degree: 450
%% duplicateId: 0
%% body:
如图所示，与水平面夹角为 $\theta$ 的绝缘斜面上固定有光滑 U 型金属导轨。质量为 $m$、电阻不可忽略的导体杆 MN 沿导轨向下运动，以大小为 $v$ 的速度进入方向垂直于导轨平面向下的匀强磁场区域，在磁场中运动一段时间 $t$ 后，速度大小变为 $2v$。运动过程中杆与导轨垂直并接触良好，导轨的电阻忽略不计，重力加速度为 $g$。杆在磁场中运动的此段时间内\xzanswer{D}

\onepicture[h=3.20cm]{figs/07.png}

\fourchoices[answer=D]
{流过杆的感应电流方向从 N 到 M}
{杆沿轨道下滑的距离为 $\frac{3}{2}vt$}
{流过杆感应电流的平均电功率等于重力的平均功率}
{杆所受安培力的冲量大小为 $mgt\sin\theta-mv$}

%% answer: D
%% memo:
\memoanswer{
A．根据右手定则，判断知流过杆的感应电流方向从 M 到 N，故 A 错误；

B．依题意，设杆切割磁感线的有效长度为 $L$，电阻为 $R$。杆在磁场中运动的此段时间内，杆受到重力、轨道支持力及沿轨道向上的安培力作用，根据牛顿第二定律可得
$mg\sin\theta-F_{\text{安}}=ma$

$F_{\text{安}}=BIL$

$I=\frac{BLv}{R}$

联立可得杆的加速度
$a=g\sin\theta-\frac{B^{2}L^{2}v}{R}$

可知，杆在磁场中运动的此段时间内做加速度逐渐减小的加速运动；若杆做匀加速直线运动，则杆运动的距离为
$s=\frac{v+2v}{2}\cdot t=\frac{3}{2}vt$

根据 $v-t$ 图像围成的面积表示位移，可知杆在时间 $t$ 内速度由 $v$ 达到 $2v$，杆真实运动的距离大于匀加速情况发生的距离，即大于 $\frac{3}{2}vt$，故 B 错误；

C．由于在磁场中运动的此段时间内，杆做加速度逐渐减小的加速运动，杆的动能增大。由动能定理可知，重力对杆所做的功大于杆克服安培力所做的功，根据 $\overline{P}=\frac{W}{t}$ 可得安培力的平均功率小于重力的平均功率，也即流过杆感应电流的平均电功率小于重力的平均功率，故 C 错误；

D．杆在磁场中运动的此段时间内，根据动量定理，可得
$mgt\sin\theta-I_{\text{安}}=m\cdot 2v-mv$

得杆所受安培力的冲量大小为 $I_{\text{安}}=mgt\sin\theta-mv$

故 D 正确。

故选 D。
}

\item
%% number: 8
%% paperName: 2023年重庆高考第 8 题 5 分
%% typeId: 2
%% type: 多选题
%% chapter: 直线运动
%% point: 运动图象
%% method: 无
%% score: 5
%% degree: 450
%% duplicateId: 0
%% body:
某实验小组测得在竖直方向飞行的无人机飞行高度 $y$ 随时间 $t$ 的变化曲线如图所示，E、F、M、N 为曲线上的点，EF、MN 段可视为两段直线，其方程分别为 $y=4t-26$ 和 $y=-2t+140$。无人机及其载物的总质量为 $2\Ukg$，取竖直向上为正方向。则\xzanswer{AB}

\onepicture[h=3.60cm]{figs/08.png}

\fourchoices[answer=AB]
{EF 段无人机的速度大小为 $4\Ums$}
{FM 段无人机的货物处于失重状态}
{FN 段无人机和装载物总动量变化量大小为 $4\Ukgms$}
{MN 段无人机机械能守恒}

%% answer: AB
%% memo:
\memoanswer{
A．根据 EF 段方程
$y=4t-26$

可知 EF 段无人机的速度大小为
$v=\frac{\Delta y}{\Delta t}=4\Ums$

故 A 正确；

B．根据 $y-t$ 图像的切线斜率表示无人机的速度，可知 FM 段无人机先向上做减速运动，后向下做加速运动，加速度方向一直向下，则无人机的货物处于失重状态，故 B 正确；

C．根据 MN 段方程
$y=-2t+140$

可知 MN 段无人机的速度为
$v'=\frac{\Delta y'}{\Delta t'}=-2\Ums$

则有
$\Delta p=mv'-mv=2\times(-2)\Ukgms-2\times4\Ukgms=-12\Ukgms$

可知 FN 段无人机和装载物总动量变化量大小为 $12\Ukgms$，故 C 错误；

D．MN 段无人机向下做匀速直线运动，动能不变，重力势能减少，无人机的机械能不守恒，故 D 错误。

故选 AB。
}

\item
%% number: 9
%% paperName: 2023年重庆高考第 9 题 5 分
%% typeId: 2
%% type: 多选题
%% chapter: 机械波
%% point: 波的图象
%% method: 无
%% score: 5
%% degree: 550
%% duplicateId: 0
%% body:
一列简谐横波在介质中沿 $x$ 轴传播，波速为 $2\Ums$，$t=0$ 时的波形图如图所示，P 为该介质中的一质点。则\xzanswer{BD}

\onepicture[h=3.00cm]{figs/09.png}

\fourchoices[answer=BD]
{该波的波长为 $14\Um$}
{该波的周期为 $8\Us$}
{$t=0$ 时质点 P 的加速度方向沿 $y$ 轴负方向}
{$0\sim2\Us$ 内质点 P 运动的路程有可能小于 $0.1\Um$}

%% answer: BD
%% memo:
\memoanswer{
A．由图可知
$\frac{3}{4}\lambda=12\Um$

解得
$\lambda=16\Um$

A 错误；

B．由
$v=\frac{\lambda}{T}$

得
$T=\frac{\lambda}{v}=\frac{16}{2}\Us=8\Us$

B 正确；

C．简谐运动的加速度总指向平衡位置，P 点位于 $y$ 轴负半轴，加速度方向沿 $y$ 轴正方向，C 错误；

D．P 点位于 $y$ 轴的负半轴，经过
$2\Us=\frac{T}{4}$

若波向 $x$ 轴负方向传播，P 向远离平衡位置方向振动，在 $0\sim2\Us$ 内质点 P 运动的路程有可能小于 $0.1\Um$，D 正确。

故选 BD。
}

\item
%% number: 10
%% paperName: 2023年重庆高考第 10 题 5 分
%% typeId: 2
%% type: 多选题
%% chapter: 曲线运动
%% point: 人造地球卫星
%% method: 无
%% score: 5
%% degree: 450
%% duplicateId: 0
%% body:
某卫星绕地心的运动视为匀速圆周运动，其周期为地球自转周期 $T$ 的 $\frac{3}{10}$，运行的轨道与地球赤道不共面（如图）。$t_{0}$ 时刻，卫星恰好经过地球赤道上 P 点正上方。地球的质量为 $M$，半径为 $R$，引力常量为 $G$。则\xzanswer{BCD}

\onepicture[h=3.20cm]{figs/10.png}

\fourchoices[answer=BCD]
{卫星距地面的高度为 $\sqrt[3]{\frac{GMT^{2}}{4\pi^{2}}}-R$}
{卫星与位于 P 点处物体的向心加速度大小比值为 $\frac{5}{9\pi R}\sqrt[3]{180\pi GMT^{2}}$}
{从 $t_{0}$ 时刻到下一次卫星经过 P 点正上方时，卫星绕地心转过的角度为 $20\pi$}
{每次经最短时间实现卫星距 P 点最近到最远的行程，卫星绕地心转过的角度比地球的多 $7\pi$}

%% answer: BCD
%% memo:
\memoanswer{
A．由题意，知卫星绕地球运转的周期为
$T'=\frac{3}{10}T$

设卫星的质量为 $m$，卫星距地面的高度为 $h$，有
$G\frac{Mm}{(R+h)^{2}}=m(R+h)(\frac{2\pi}{T'})^{2}$

联立，可求得
$h=(\frac{9GMT^{2}}{400\pi^{2}})^{\frac{1}{3}}-R$

故 A 错误；

B．卫星的向心加速度大小
$a_{1}=(R+h)\omega'^{2}=(R+h)(\frac{2\pi}{T'})^{2}$

位于 P 点处物体的向心加速度大小
$a_{2}=R\omega^{2}=R(\frac{2\pi}{T})^{2}$

可得
$\frac{a_{1}}{a_{2}}=\frac{R+h}{R}(\frac{T}{T'})^{2}=\frac{5}{9\pi R}(180\pi GMT^{2})^{\frac{1}{3}}$

故 B 正确；

CD．要想卫星再次在 P 点的正上方，则只能是题中两个轨道的交点，因此要实现出现在正上方，第一种情形是经过一段时间都回到了当前点，即各自转动整数圈，最小公倍数为 3，此时卫星转动 10 圈，即转动角度为 $20\pi$，第二种情形是都转动整数圈加半圈，此时最小公倍数为 1.5，卫星转动五圈，此时相距最远，转动角度相差 $7\pi$，故 CD 正确。

故选 BCD。
}

\item
%% number: 11
%% paperName: 2023年重庆高考第 11 题 8 分
%% typeId: 4
%% type: 实验题
%% chapter: 机械振动
%% point: 单摆测重力加速度
%% method: 无
%% score: 8
%% degree: 550
%% duplicateId: 0
%% body:
某实验小组用单摆测量重力加速度。所用实验器材有摆球、长度可调的轻质摆线、刻度尺、50 分度的游标卡尺、摄像装置等。

\twopicture[numA=甲, numB=乙, h=4.20cm]
{figs/11a.png}
{figs/11b.png}

\begin{enumerate}
	\item
	用游标卡尺测量摆球直径 $d$。当量爪并拢时，游标尺和主尺的零刻度线对齐。放置摆球后游标卡尺示数如图甲所示，则摆球的直径 $d$ 为\tkanswer[1.60cm]{19.20} $\Umm$。
	\item
	用摆线和摆球组成单摆，如图乙所示。当摆线长度 $l=990.1 \Umm$ 时，记录并分析单摆的振动视频，得到单摆的振动周期 $T=2.00 \Us$，由此算得重力加速度 $g$ 为\tkanswer[1.60cm]{9.86} $\Umsq$（保留 3 位有效数字）。
	\item
	改变摆线长度 $l$，记录并分析单摆的振动视频，得到相应的振动周期。他们发现，分别用 $l$ 和 $l+\frac{d}{2}$ 作为摆长，这两种计算方法得到的重力加速度数值的差异大小 $\Delta g$ 随摆线长度 $l$ 的变化曲线如图所示。由图可知，该实验中，随着摆线长度 $l$ 的增加，$\Delta g$ 的变化特点是\tkanswer[2.40cm]{逐渐减小}，原因是\tkanswer[5.00cm]{随着摆线长度 $l$ 的增加，则 $l+\frac{d}{2}$ 越接近于 $l$，此时计算得到的 $g$ 的差值越小}。
\end{enumerate}

\onepicture[h=4.00cm, align=right]{figs/11c.png}

%% answer:
\jdanswer{
\begin{enumerate}
	\item
	19.20

	\item
	9.86

	\item
	随着摆线长度 $l$ 的增加，$\Delta g$ 逐渐减小；原因是随着摆线长度 $l$ 的增加，则 $l+\frac{d}{2}$ 越接近于 $l$，此时计算得到的 $g$ 的差值越小
\end{enumerate}
}
%% memo:
\memoanswer{
（1）用游标卡尺测量摆球直径 $d=19\Umm+0.02\Umm\times10=19.20\Umm$

（2）单摆的摆长为 $L=990.1\Umm+\frac{1}{2}\times19.20\Umm=999.7\Umm$

根据 $T=2\pi\sqrt{\frac{L}{g}}$，可得 $g=\frac{4\pi^{2}L}{T^{2}}$，带入数据
$g=\frac{4\times3.14^{2}\times0.9997}{2^{2}}\Umsq=9.86\Umsq$

（3）由图可知，随着摆线长度 $l$ 的增加，$\Delta g$ 逐渐减小，原因是随着摆线长度 $l$ 的增加，则 $l+\frac{d}{2}$ 越接近于 $l$，此时计算得到的 $g$ 的差值越小。
}

\item
%% number: 12
%% paperName: 2023年重庆高考第 12 题 8 分
%% typeId: 4
%% type: 实验题
%% chapter: 交流电
%% point: 变压器
%% method: 无
%% score: 8
%% degree: 550
%% duplicateId: 0
%% body:
一兴趣小组拟研究某变压器的输入和输出电压之比，以及交流电频率对输出电压的影响。题图 1 为实验电路图，其中 L$_{1}$ 和 L$_{2}$ 为变压器的原、副线圈，S$_{1}$ 和 S$_{2}$ 为开关，P 为滑动变阻器 $R_{P}$ 的滑片，$R$ 为电阻箱，$E$ 为正弦式交流电源（能输出电压峰值不变、频率可调的交流电）。
\begin{enumerate}
	\item
	闭合 S$_{1}$，用多用电表交流电压挡测量线圈 L$_{1}$ 两端的电压。滑片 P 向右滑动后，与滑动前相比，电表的示数\tkanswer[1.60cm]{变大}（选填“变大”、“不变”或“变小”）。
	\item
	保持 S$_{2}$ 断开状态，调整 $E$ 输出的交流电频率为 $50\UHz$，滑动滑片 P，用多用电表交流电压挡测得线圈 L$_{1}$ 两端的电压为 $2\,500\UmV$ 时，用示波器测得线圈 L$_{2}$ 两端电压 $u$ 随时间 $t$ 的变化曲线如图所示，则线圈 L$_{1}$ 两端与 L$_{2}$ 两端的电压比值为\tkanswer[1.60cm]{12.6}（保留 3 位有效数字）。
	\item
	闭合 S$_{2}$，滑动 P 到某一位置并保持不变。分别在 $E$ 输出的交流电频率为 $50\UHz$、$1\,000\UHz$ 的条件下，改变 $R$ 的阻值，用多用电表交流电压挡测量线圈 L$_{2}$ 两端的电压 $U$，得到 $U-R$ 关系曲线如图 3 所示。用一个阻值恒为 $20\UO$ 的负载 $R_{0}$ 替换电阻箱 $R$，由图可知，当频率为 $1\,000\UHz$ 时，$R_{0}$ 两端的电压为\tkanswer[1.80cm]{272} $\UmV$；当频率为 $50\UHz$ 时，为保持 $R_{0}$ 两端的电压不变，需要将 $R_{0}$ 与一个阻值为\tkanswer[1.60cm]{7.94} $\UO$ 的电阻串联。（均保留 3 位有效数字）
\end{enumerate}

\threepicture[numA=1, numB=2, numC=3, hA=3.00cm, hB=2.70cm, hC=4.40cm, align=right]
{figs/12a.png}
{figs/12b.png}
{figs/12c.png}

%% answer:
\jdanswer{
\begin{enumerate}
	\item
	变大

	\item
	12.6

	\item
	272，7.94
\end{enumerate}
}
%% memo:
\memoanswer{
（1）闭合 S$_{1}$，滑动变阻器 $R_{p}$ 是分压接法，滑片 P 向右滑动后，用多用电表交流电压挡测量线圈 L$_{1}$ 两端的电压。线圈 L$_{1}$ 两端的电压增大，因此与滑动前相比，电表的示数变大。

（2）保持 S$_{2}$ 断开状态，调整 $E$ 输出的交流电频率为 $50\UHz$，多用电表交流电压挡测得线圈 L$_{1}$ 两端的电压为 $U_{1}=2500\UmV$。线圈 L$_{2}$ 两端电压 $u$ 随时间 $t$ 的变化曲线如图所示，由 $u-t$ 图像可得，线圈 L$_{2}$ 两端电压为
$U_{2}=\frac{U_{m}}{\sqrt{2}}=\frac{280}{\sqrt{2}}\UmV$

则线圈 L$_{1}$ 两端与 L$_{2}$ 两端的电压比值为
$\frac{U_{1}}{U_{2}}=\frac{2500}{\frac{280}{\sqrt{2}}}=12.6$

（3）闭合 S$_{2}$，滑动 P 到某一位置并保持不变。由 $U-R$ 关系曲线可得，当频率为 $1000\UHz$ 时，当负载电阻 $R_{0}=20\UO$ 时，$R_{0}$ 两端的电压为 $U_{R0}=272\UmV$。

当频率为 $50\UHz$ 时，由 $U-R$ 关系曲线可得线圈 L$_{2}$ 两端的电压为 $U'=380\UmV$。要保持 $R_{0}$ 两端的电压不变，需给 $R_{0}$ 串联一电阻，此串联电阻值为
$R_{\text{串}}=\frac{U'-U_{R0}}{\frac{U_{R0}}{R_{0}}}=\frac{380-272}{\frac{272}{20}}\UO=7.94\UO$
}

\item
%% number: 13
%% paperName: 2023年重庆高考第 13 题 11 分
%% typeId: 6
%% type: 计算题
%% chapter: 功和能
%% point: 动能定理
%% method: 无
%% score: 11
%% degree: 750
%% duplicateId: 0
%% body:
机械臂广泛应用于机械装配。若某质量为 $m$ 的工件（视为质点）被机械臂抓取后，在竖直平面内由静止开始斜向上做加速度大小为 $a$ 的匀加速直线运动，运动方向与竖直方向夹角为 $\theta$，提升高度为 $h$，如图所示。求：
\begin{enumerate}
	\item
	提升高度为 $h$ 时，工件的速度大小；
	\item
	在此过程中，工件运动的时间及合力对工件做的功。
\end{enumerate}

\onepicture[h=3.80cm, align=right]{figs/13.png}

%% answer:
\jdanswer{
\begin{enumerate}
	\item
	$\sqrt{\frac{2ah}{\cos\theta}}$

	\item
	$\sqrt{\frac{2h}{a\cos\theta}}$，$\frac{mah}{\cos\theta}$
\end{enumerate}
}
%% memo:
\memoanswer{
（1）根据匀变速直线运动位移与速度关系有
$v_{0}^{2}=2a\frac{h}{\cos\theta}$

解得 $v_{0}=\sqrt{\frac{2ah}{\cos\theta}}$

（2）根据速度公式有 $v_{0}=at$

解得 $t=\sqrt{\frac{2h}{a\cos\theta}}$

根据动能定理有 $W_{\text{合}}=\frac{1}{2}mv_{0}^{2}$

解得 $W_{\text{合}}=\frac{mah}{\cos\theta}$
}

\item
%% number: 14
%% paperName: 2023年重庆高考第 14 题 12 分
%% typeId: 6
%% type: 计算题
%% chapter: 曲线运动
%% point: 圆周运动的应用
%% method: 无
%% score: 12
%% degree: 650
%% duplicateId: 0
%% body:
如图所示，桌面上固定有一半径为 $R$ 的水平光滑圆轨道，M、N 为轨道上的两点，且位于同一直径上，P 为 MN 段的中点。在 P 点处有一加速器（大小可忽略），小球每次经过 P 点后，其速度大小都增加 $v_{0}$。质量为 $m$ 的小球 1 从 N 处以初速度 $v_{0}$ 沿轨道逆时针运动，与静止在 M 处的小球 2 发生第一次弹性碰撞，碰后瞬间两球速度大小相等。忽略每次碰撞时间。求：
\begin{enumerate}
	\item
	球 1 第一次经过 P 点后瞬间向心力的大小；
	\item
	球 2 的质量；
	\item
	两球从第一次碰撞到第二次碰撞所用时间。
\end{enumerate}

\onepicture[h=3.80cm, align=right]{figs/14.png}

%% answer:
\jdanswer{
\begin{enumerate}
	\item
	$\frac{4mv_{0}^{2}}{R}$

	\item
	$3m$

	\item
	$\frac{5\pi R}{6v_{0}}$
\end{enumerate}
}
%% memo:
\memoanswer{
（1）球 1 第一次经过 P 点后瞬间速度变为 $2v_{0}$，所以
$F_{n}=m\frac{(2v_{0})^{2}}{R}=4m\frac{v_{0}^{2}}{R}$

（2）球 1 与球 2 发生弹性碰撞，且碰后速度大小相等，说明球 1 碰后反弹，则
$m\cdot2v_{0}=-mv+m'v$

$\frac{1}{2}m(2v_{0})^{2}=\frac{1}{2}mv^{2}+\frac{1}{2}m'v^{2}$

联立解得
$v=v_{0}$，$m'=3m$

（3）设两球从第一次碰撞到第二次碰撞所用时间为 $\Delta t$，则
$t_{1}=\frac{\pi R}{2v_{0}}$

$v_{0}t_{2}+2v_{0}t_{2}=\pi R$

所以 $\Delta t=t_{1}+t_{2}=\frac{5\pi R}{6v_{0}}$
}

\item
%% number: 15
%% paperName: 2023年重庆高考第 15 题 18 分
%% typeId: 6
%% type: 计算题
%% chapter: 磁场
%% point: 带电粒子在组合场中的运动
%% method: 无
%% score: 18
%% degree: 350
%% duplicateId: 0
%% body:
某同学设计了一种粒子加速器的理想模型。如图所示，$xOy$ 平面内，$x$ 轴下方充满垂直于纸面向外的匀强磁场，$x$ 轴上方被某边界分割成两部分，一部分充满匀强电场（电场强度与 $y$ 轴负方向成 $\alpha$ 角），另一部分无电场，该边界与 $y$ 轴交于 M 点，与 $x$ 轴交于 N 点。只有经电场到达 N 点、与 $x$ 轴正方向成 $\alpha$ 角斜向下运动的带电粒子才能进入磁场。从 M 点向电场内发射一个比荷为 $\frac{q}{m}$ 的带电粒子 A，其速度大小为 $v_{0}$、方向与电场方向垂直，仅在电场中运动时间 $T$ 后进入磁场，且通过 N 点的速度大小为 $2v_{0}$。忽略边界效应，不计粒子重力。
\begin{enumerate}
	\item
	求角度 $\alpha$ 及 M、N 两点的电势差。
	\item
	在该边界上任意位置沿与电场垂直方向直接射入电场内的、比荷为 $\frac{q}{m}$ 的带电粒子，只要速度大小适当，就能通过 N 点进入磁场，求 N 点横坐标及此边界方程。
	\item
	若粒子 A 第一次在磁场中运动时磁感应强度大小为 $B_{1}$，以后每次在磁场中运动时磁感应强度大小为上一次的一半，则粒子 A 从 M 点发射后，每次加速均能通过 N 点进入磁场。求磁感应强度大小 $B_{1}$ 及粒子 A 从发射到第 $n$ 次通过 N 点的时间。
\end{enumerate}

\onepicture[h=3.80cm, align=right]{figs/15.png}

%% answer:
\jdanswer{
\begin{enumerate}
	\item
	$30^\circ$，$\frac{3mv_{0}^{2}}{2q}$

	\item
	$\frac{3\sqrt{3}v_{0}T}{4}$，$y=\frac{v_{0}T}{4}-\frac{\sqrt{3}}{9}x$

	\item
	$B_{1}=\frac{\sqrt{3}m}{6qT}$，$t=(3+\frac{10\sqrt{3}\pi}{3})2^{n-1}T+(2+\frac{10\sqrt{3}\pi}{3})T$
\end{enumerate}
}
%% memo:
\memoanswer{
（1）粒子 M 点垂直于电场方向入射，粒子在电场中做类平抛运动，沿电场方向做匀加速直线运动，垂直于电场方向做匀速直线运动，在 N 点将速度沿电场方向与垂直于电场方向分解，在垂直于电场方向上有
$2v_{0}\cos2\alpha=v_{0}$

解得 $\alpha=30^\circ$

粒子从 $M\to N$ 过程，根据动能定理有
$qU_{MN}=\frac{1}{2}m(2v_{0})^{2}-\frac{1}{2}mv_{0}^{2}$

解得 $U_{MN}=\frac{3mv_{0}^{2}}{2q}$

（2）对于从 M 点射入的粒子，沿初速度方向的位移为
$x_{0}=v_{0}T$

沿电场方向的位移为
$y_{0}=\frac{2v_{0}\sin2\alpha}{2}T$

令 N 点横坐标为 $x_{N}$，根据几何关系有
$x_{N}=x_{0}\cos\alpha+y_{0}\sin\alpha$

解得 $x_{N}=\frac{3\sqrt{3}v_{0}T}{4}$

根据上述与题意可知，令粒子入射速度为 $v$，则通过 N 点进入磁场的速度为 $2v$，令边界上点的坐标为 $(x,y)$，则在沿初速度方向上有
$(x_{N}-x-y\tan\alpha)\cos\alpha=vt$

在沿电场方向有
$\frac{y}{\cos\alpha}+(x_{N}-x-y\tan\alpha)\sin\alpha=\frac{2v\sin2\alpha}{2}t$

解得 $y=\frac{v_{0}T}{4}-\frac{\sqrt{3}}{9}x$

（3）由上述结果可知电场强度 $E=\frac{U_{MN}}{y_{0}}$

解得 $E=\frac{\sqrt{3}mv_{0}}{qT}$

设粒子 A 第 $n$ 次在磁场中做圆周运动的线速度为 $v_{n}$，可得第 $n+1$ 次在 N 点进入磁场的速度为
$v_{n+1}=\frac{v_{n}}{\cos2\alpha}=2v_{n}$

第一次在 N 点进入磁场的速度大小为 $2v_{0}$，可得
$v_{n}=2^{n}v_{0},(n=1,2,3\cdots)$

设粒子 A 第 $n$ 次在磁场中运动时的磁感应强度为 $B_{n}$，由题意可得
$B_{n}=\frac{B_{1}}{2^{n-1}},(n=1,2,3\cdots)$

由洛伦兹力提供向心力得
$qv_{n}B_{n}=m\frac{v_{n}^{2}}{r_{n}}$

联立解得
$r_{n}=\frac{4^{n}mv_{0}}{2qB_{1}}$

粒子 A 第 $n$ 次在磁场中的运动轨迹如图所示

\onepicture[h=4.60cm]{figs/15_an.png}

粒子每次在磁场中运动轨迹的圆心角均为 $300^\circ$，第 $n$ 次离开磁场的位置 C 与 N 的距离等于 $r_{n}$，由 C 到 N 由动能定理得
$qEr_{n}\sin30^\circ=\frac{1}{2}mv_{n+1}^{2}-\frac{1}{2}mv_{n}^{2}$

联立上式解得 $B_{1}=\frac{\sqrt{3}m}{6qT}$

由类平抛运动沿电场方向的运动可得，粒子 A 第 $n$ 次在电场中运动的时间为
$t_{1n}=\frac{2^{n}v_{0}\sin60^\circ}{\frac{qE}{m}}=2^{n-1}T$

粒子 A 第 $n$ 次在磁场中运动的周期为
$T'=\frac{2\pi r_{n}}{v_{n}}=2^{n+1}\sqrt{3}\pi T$

粒子 A 第 $n$ 次在磁场中运动的时间为
$t_{2n}=\frac{300^\circ}{360^\circ}T'=\frac{5}{6}\cdot2^{n+1}\sqrt{3}\pi T$

设粒子 A 第 $n$ 次在电场边界 MN 与 $x$ 轴之间的无场区域的位移为 $x_{n}$，边界与 $x$ 轴负方向的夹角为 $\beta$，则根据边界方程可得
$\tan\beta=\frac{\sqrt{3}}{9}$，$\sin\beta=\frac{1}{2\sqrt{7}}$

由正弦定理可得
$\frac{x_{n}}{\sin\beta}=\frac{r_{n}}{\sin(180^\circ-30^\circ-\beta)}$

解得
$x_{n}=\frac{4^{n}v_{0}T}{2}$

粒子 A 第 $n$ 次在电场边界 MN 与 $x$ 轴之间运动的时间为
$t_{3n}=\frac{x_{n}}{v_{n}}=2^{n-1}T$

粒子 A 从发射到第 $n$ 次通过 N 点的过程，在电场中运动 $n$ 次，在磁场和无场区域中均运动 $n-1$ 次，则所求时间
$t=(2^{0}+2^{1}+\cdots+2^{n-1})T+\frac{5}{6}\cdot\sqrt{3}\pi T(2^{2}+2^{3}+\cdots+2^{n})+(2^{0}+2^{1}+\cdots+2^{n-2})T$

由等比数列求和得
$t=(2^{n}-1)T+\frac{5}{6}\cdot\sqrt{3}\pi T[4(2^{n-1}-1)]+(2^{n-1}-1)T$

解得
$t=(3+\frac{10\sqrt{3}\pi}{3})2^{n-1}T+(2+\frac{10\sqrt{3}\pi}{3})T$

\textbf{永动机模型错误！}

（说明：本题模型使粒子每经一次循环能量增加，构成永动机模型，物理上不能实现；以上为原卷参考答案。）
}

\end{enumerate}

\end{document}
